\chapter{Validação, comparação externa e uso das evidências}
\label{ch:validacao}

\section{Estratégia de validação}
\label{sec:estrategia-validacao}

A validação responde a duas perguntas. Primeiro: o sistema executa as regras de cálculo como foram definidas? Segundo: os dados disponíveis são suficientes para sustentar a interpretação do resultado? As duas perguntas são diferentes. Um cálculo pode estar correto e, ainda assim, usar uma estimativa de grupo ou uma distância aproximada. Esses casos exigem cautela.

Por isso, o \CabotageLens{} verifica cada etapa separadamente. Primeiro, escolhem-se dois portos: um para iniciar a navegação analisada e outro para encerrá-la. O sistema lê a lista de escalas em ordem de data dentro de uma única viagem. O recorte é aceito quando a lista contém o porto inicial e, mais adiante, o porto final. Ele pode representar uma navegação direta ou incluir paradas intermediárias.

Depois, o sistema verifica a carga a bordo e a distância entre cada par de escalas consecutivas. Cada uma dessas pernas recebe o nome de \emph{subtrecho}. O consumo da viagem deve ser a soma dos consumos de todos os subtrechos entre os dois portos escolhidos. A fonte do valor de consumo por tonelada e por milha também deve ficar registrada. Por fim, a sequência de portos usada como rota deve pertencer a uma única viagem. O sistema não pode unir trechos de viagens diferentes.

Um registro real ajuda a distinguir as etapas. Na viagem \codecell{voyage\_9612791\_00011}, o navio de IMO 9612791 passou por Santos, Suape, Pecém e Manaus, nessa ordem. O recorte Santos--Manaus contém Santos--Suape, Suape--Pecém e Pecém--Manaus. A validação confere separadamente a carga a bordo, a distância, o trabalho de transporte e o combustível de cada subtrecho. Também confere que a intensidade de \(7{,}43~\mathrm{g/(t\,nm)}\) veio diretamente do IMO no EU MRV de 2023. Assim, o sistema mantém os dados necessários para explicar o número final.

\begin{table}[htbp]
\centering
\small
\caption{Etapas verificadas na validação do sistema.}
\label{tab:camadas-validacao}
\begin{tabularx}{\textwidth}{L{0.27\textwidth}Y Y}
\toprule
Camada & O que é verificado & O que deve ficar registrado \\
\midrule
Integridade da ANTAQ & Ordem das escalas, identidade da viagem, movimentos de carga e nomes equivalentes de portos. & Viagens e subtrechos coerentes para reconstruir a atividade.\\
Associação ANTAQ--MRV & Correspondência por IMO, período positivo mais recente e estimativas feitas com navios semelhantes. & Origem do indicador de consumo usado em cada viagem.\\
Cálculo marítimo & Carga, distância, produto carga--distância e combustível em cada perna. & Recorte histórico de cada viagem e valor representativo entre os dois portos.\\
Seleção da rota & Continuidade na mesma viagem, preferência pela ligação direta e menor sequência completa. & Sequência e distância da rota separadas do conjunto de viagens usado na estatística.\\
Testes automatizados & Casos diretos, múltiplas escalas, sequências de portos diferentes, valores extremos, pesos nulos e dados incompletos. & Resposta esperada e regra executada pelo código.\\
Comparação externa & Modo apontado como menos emissor, ordem de grandeza, itens incluídos e premissas. & Classificação compatível com o limite da comparação.\\
\bottomrule
\end{tabularx}
\end{table}

\section{Integridade e cobertura da base marítima}
\label{sec:validacao-base-maritima}

Os indicadores desta seção são resultados do processamento da base ANTAQ de 2025. Eles não são exemplos hipotéticos. A tabela de Atracação forneceu \num{7103} ocorrências de entrada de navio em porto. Cada ocorrência é uma \emph{chamada portuária}. O sistema ordenou essas chamadas por navio e por data. Chamadas consecutivas do mesmo navio no mesmo porto foram reunidas em uma única posição da sequência. Essa posição recebe o nome de \emph{parada}.

As chamadas foram organizadas em \num{6797} paradas distribuídas por \num{1324} viagens reconstruídas. Duas paradas consecutivas formam uma perna navegada, chamada de \emph{subtrecho}. Na parte Santos--Suape--Pecém--Manaus da viagem \codecell{voyage\_9612791\_00011}, as quatro paradas formam três subtrechos: Santos--Suape, Suape--Pecém e Pecém--Manaus.

Também foram encontrados \num{389} números IMO distintos. O IMO funciona como a identificação internacional do navio. Entre esses IMOs, \num{243} possuem um indicador positivo de consumo por trabalho realizado no catálogo MRV.

No nível das viagens, \num{788} receberam um valor específico porque o IMO foi encontrado no MRV. Nas outras \num{536}, o sistema usou um valor calculado para navios do mesmo tipo. Essa estimativa de grupo é chamada de \emph{fallback}. Ela permite calcular a viagem, mas não é um valor específico daquele navio.

O resultado registra qual das duas fontes foi usada. Esse registro da origem do dado recebe o nome de \emph{proveniência}. Assim, o leitor consegue distinguir um valor ligado diretamente ao IMO de uma estimativa baseada em outros navios.

A geração da matriz começou com \num{5473} subtrechos candidatos. Esse total decorre das \num{6797} paradas: uma viagem com \(n\) paradas tem \(n-1\) subtrechos. Ao somar todas as viagens, a conta é \(6797-1324=5473\).

O sistema padronizou nomes diferentes usados para o mesmo porto. Também reuniu \num{35} transições consecutivas que, depois da padronização, apontavam para o mesmo porto. Restaram \num{5438} subtrechos calculados. Desses, \num{5023} usam distância da matriz marítima. Os outros \num{415} usam uma aproximação de grande círculo, identificada como haversine. Na intensidade, \num{3290} subtrechos usam IMO e \num{2148} usam fallback robusto.

Uma viagem pode produzir vários recortes entre portos. Dentro da viagem \codecell{voyage\_9612791\_00011}, o prefixo Santos--Suape--Pecém--Manaus gera seis recortes nesse sentido: Santos--Suape, Santos--Pecém, Santos--Manaus, Suape--Pecém, Suape--Manaus e Pecém--Manaus. A viagem completa ainda retorna de Manaus a Santos e pode fornecer outras ligações. Os três recortes entre portos consecutivos são diretos. Os outros três contêm uma ou duas paradas intermediárias. Assim, \num{1324} viagens podem gerar muito mais que \num{1324} recortes.

Cada recorte permanece dentro de uma única viagem. Ele começa em uma parada escolhida como partida e termina em outra parada visitada mais tarde como chegada. O nome técnico desse recorte é \emph{recorte histórico viagem--OD}, em que OD significa origem--destino. O processamento produz um arquivo final usado pelo aplicativo. Esse arquivo contém \num{12025} recortes históricos viagem--OD.

Um recorte fica pronto para entrar no cálculo quando passa por quatro verificações. Primeiro, a parte aproveitada da viagem deve seguir a direção pedida. No cálculo Santos--Manaus, ela começa na escala de Santos e termina na escala de Manaus; uma sequência Manaus--Suape--Santos não serve. Segundo, nenhum dos trechos navegados entre esses portos pode estar ausente. Terceiro, deve existir uma intensidade obtida pelo IMO ou estimada e identificada como fallback. Quarto, o trabalho de transporte deve ser maior que zero para o recorte receber peso na mediana; esse trabalho é a soma da carga a bordo multiplicada pela distância de cada trecho. Um recorte com resultado igual a zero permanece registrado, mas não altera a mediana quando existem resultados positivos. Se todos forem iguais a zero, o sistema usa a mediana sem pesos e registra a exceção.

Os \num{12025} recortes cobrem \num{177} ligações com sentido definido. Por exemplo, Santos--Manaus e Manaus--Santos são ligações diferentes. Se o navio visitar os portos escolhidos mais de uma vez durante a mesma viagem, essa viagem continua fornecendo somente um recorte para a ligação. Isso pode ser verificado na viagem real \codecell{voyage\_9697002\_00002}, cuja sequência canônica foi Santos--Itapoá--Rio de Janeiro--Suape--Pecém--Manaus--Suape--Santos. Para Suape--Santos, a viagem oferece uma passagem longa a partir da primeira escala em Suape e uma passagem direta a partir da segunda. A regra conserva somente a passagem direta. Se nenhuma passagem direta existir dentro da viagem, usa a parte completa de menor distância entre os dois portos.

\begin{table}[htbp]
\centering
\small
\caption{Indicadores de integridade e cobertura do arquivo marítimo processado.}
\label{tab:integridade-maritima}
\begin{tabularx}{\textwidth}{Y r Y}
\toprule
Indicador & Valor & Como ler o indicador \\
\midrule
Viagens ANTAQ & \num{1324} & Sequências independentes reconstruídas a partir das chamadas.\\
IMOs únicos / IMOs associados ao MRV & \num{389} / \num{243} & Cobertura de intensidade específica por embarcação.\\
Viagens com IMO exato / fallback & \num{788} / \num{536} & Fonte da intensidade atribuída à viagem.\\
Subtrechos calculados & \num{5438} & Pernas consecutivas com porto e distância resolvidos.\\
Recortes históricos viagem--OD & \num{12025} & Partes de viagens entre duas paradas escolhidas; uma viagem pode gerar vários recortes.\\
Ligações portuárias com sentido definido & \num{177} & Relações disponíveis ao cálculo multimodal.\\
\bottomrule
\end{tabularx}
\end{table}

\section{Testes das regras de cálculo}
\label{sec:testes-regras-maritimas}

Os testes unitários combinam dois tipos de verificação. Casos sintéticos pequenos isolam uma regra e permitem calcular a resposta à mão. Um teste de regressão usa os movimentos reais da viagem \codecell{voyage\_9612791\_00011} para conferir a reconstrução de Santos--Suape--Pecém--Manaus e os registros de auditoria. Os casos sintéticos não são apresentados como resultados da ANTAQ; eles funcionam apenas como instrumentos de teste.

No teste de viagem direta, há somente uma perna entre os dois portos. O teste verifica se o sistema encontra o registro positivo mais recente do IMO no MRV. No teste real com múltiplas escalas, o consumo esperado é a soma de Santos--Suape, Suape--Pecém e Pecém--Manaus: \(120.302{,}670+63.078{,}304+108.578{,}413=291.959{,}387\)~kg para a atividade histórica total do navio no recorte Santos--Manaus. Dessa forma, o teste detecta se o código usar apenas uma distância direta ou uma única carga em toda a viagem.

Outro teste escolhe os mesmos portos de partida e chegada em duas viagens controladas: uma direta e outra com escala intermediária. Cada sequência completa de portos recebe o nome de \emph{corredor}. As duas viagens devem ajudar a calcular o indicador típico de consumo da ligação. Porém, apenas um corredor completo fornece a sequência e a distância usadas no cenário. O teste impede a montagem de uma rota artificial com pernas de viagens diferentes. Na base real, a mesma separação ocorre entre a ligação direta de \codecell{voyage\_9612789\_00004} e os corredores com escalas, como \codecell{voyage\_9612791\_00011}.

A regra para valores extremos também é testada. Um valor ligado diretamente a um IMO é preservado. Para estimar um tipo de navio, o sistema retira 1\% dos menores valores e 1\% dos maiores quando a amostra permite. Depois, calcula a média do restante. Essa regra recebe o nome de média aparada. Em amostras pequenas, usa-se a mediana. Para a classe, usa-se a média aparada gravada no arquivo ou, se ela não existir, a mediana. O teste confere o valor e a fonte registrada.

Há ainda testes para carga acumulada não negativa e peso zero. Outros testes verificam nomes diferentes do mesmo terminal e a distância aproximada por haversine. O sistema também exige o formato marítimo 3 para impedir a leitura silenciosa de um arquivo incompatível. Se todos os recortes de uma ligação tiverem trabalho de transporte igual a zero, o sistema aplica a mediana simples das intensidades e registra essa exceção.

\section{Auditoria da ligação Santos--Manaus}
\label{sec:auditoria-santos-manaus}

Esta auditoria examina como o indicador histórico de Santos--Manaus foi construído. Ela não é, ainda, uma execução com a carga informada por um usuário. O sistema escolhe Santos como porto de partida e Manaus como porto de chegada. Em seguida, lê cronologicamente as escalas de cada viagem histórica. Entram os recortes que começam em uma parada em Santos e terminam em uma parada posterior em Manaus.

Foram encontrados \num{89} recortes nos quais o navio saiu de Santos e chegou depois a Manaus na mesma viagem, sem faltar nenhum subtrecho intermediário e com os dados necessários para calcular a intensidade e o peso estatístico. Eles estão distribuídos em \num{22} sequências completas de portos, ou corredores. Há um recorte direto. Há também corredores com escalas. Exemplos observados incluem Santos--Suape--Pecém--Manaus em \codecell{voyage\_9612791\_00011}; Santos--Rio de Janeiro--Suape--Pecém--Manaus em \codecell{voyage\_9784661\_00001}; e Santos--Navegantes--Pecém--Manaus em \codecell{voyage\_9852365\_00011}. O último não passa por Suape.

Em cada recorte histórico, o sistema multiplica a carga que estava a bordo pela distância navegada em cada subtrecho. Depois, soma esses produtos. O total mede quanto peso foi transportado e por qual distância. Sua unidade é tonelada-milha náutica, ou t$\cdot$nm. Esse total recebe o nome de \emph{trabalho de transporte}.

Os 89 recortes formam o conjunto usado para calcular o indicador típico de consumo de Santos--Manaus. Destes, \num{40} usam valor específico por IMO e \num{49} usam fallback de tipo. Embora as contagens sejam próximas, os fallbacks representam \num{59.54}\% do trabalho de transporte. Esse percentual é importante porque um recorte com maior soma de carga a bordo multiplicada pela distância de cada trecho recebe mais peso.

Cada recorte usa o indicador do navio que realizou a viagem, expresso em gramas de combustível por tonelada transportada e por milha náutica. Esse indicador recebe o nome de \emph{intensidade}. Para obter um único valor para Santos--Manaus, o sistema ordena as intensidades e acumula o trabalho de transporte. Ele escolhe a primeira intensidade em que o total acumulado alcança 50\%. Essa regra recebe o nome de \emph{mediana ponderada pelo trabalho de transporte}.

O valor aplicado é \num{9.322050}~g/(t$\cdot$nm). A média aritmética ponderada dos mesmos recortes históricos é \num{25.243618}~g/(t$\cdot$nm). Ela é mantida como diagnóstico, mas não entra no cenário. A diferença mostra que intensidades altas exercem forte influência sobre a média.

O ponto de corte também pode ser reproduzido com registros individuais. Os 89 recortes somam \(3.153.328.821{,}755~\mathrm{t\,nm}\), de modo que 50\% correspondem a \(1.576.664.410{,}877~\mathrm{t\,nm}\). Antes de incluir \codecell{voyage\_9697002\_00002}, o acumulado estava em 49,168\%. O recorte dessa viagem seguiu Santos--Itapoá--Rio de Janeiro--Suape--Pecém--Manaus e acrescentou \(34.357.307{,}013~\mathrm{t\,nm}\), elevando o acumulado para 50,257\%. Por isso, sua intensidade de \(9{,}322050~\mathrm{g/(t\,nm)}\) define a mediana ponderada. A fonte registrada foi a média aparada em 1\% do tipo-padrão documentado \emph{container ship}, usado porque não havia correspondência individual aplicável pelo IMO nem classe ou tipo específico informado.

Depois de construir o indicador com as viagens históricas, o sistema escolhe um único corredor para a execução do usuário. Primeiro, procura um recorte direto entre os dois portos. Se não encontrar, procura o corredor completo de menor distância. O corredor selecionado fornece os portos e a distância modelada do cálculo. O código chama esse conjunto de \emph{caminho}.

Como existe um recorte direto observado, o caminho escolhido é Santos--Manaus. Esse recorte pertence a \codecell{voyage\_9612789\_00004}, do navio de IMO 9612789. A matriz marítima atribui \num{6112}~km, equivalentes a \num{3300.216}~nm, ao subtrecho, e a carga histórica reconstruída foi \(11.584{,}165\)~t. A distância não é uma trajetória medida por posições do navio. Os outros 21 corredores continuam no conjunto histórico usado para calcular a intensidade. Portanto, todos ajudam a definir o indicador de consumo, mas somente o caminho direto fornece a distância aplicada à remessa do usuário.

A carga a bordo reconstruída nos 89 recortes históricos tem uma função estatística: ela define o trabalho de transporte e o peso de cada recorte. Essa carga não é usada como a massa da nova remessa. No teste de execução registrado neste trabalho, o sistema recebe uma única entrada controlada de 14~t e aplica a ela a intensidade de \num{9.322050}~g/(t$\cdot$nm) e a distância selecionada. Assim, o histórico forma o indicador; a entrada de teste forma um cenário único e não é apresentada como carga observada na ANTAQ.

\begin{table}[htbp]
\centering
\small
\caption{Auditoria do estimador marítimo para Santos--Manaus.}
\label{tab:auditoria-santos-manaus}
\begin{tabularx}{\textwidth}{Y r Y}
\toprule
Indicador & Valor & Como ler o indicador \\
\midrule
Recortes históricos resolvidos & \num{89} & Partes de viagens que começam em Santos e terminam mais tarde em Manaus.\\
Corredores observados & \num{22} & Uma ligação direta e sequências completas com escalas.\\
Fonte por IMO / fallback de tipo & \num{40} / \num{49} & Composição da origem das intensidades dos recortes.\\
Trabalho associado a fallback & \num{59.54}\% & Participação do fallback no peso total.\\
Mediana ponderada & \num{9.322050} g/(t$\cdot$nm) & Intensidade aplicada à ligação Santos--Manaus.\\
Média ponderada diagnóstica & \num{25.243618} g/(t$\cdot$nm) & Valor auxiliar, não aplicado ao cenário.\\
Rota selecionada & Santos--Manaus & Sequência direta observada usada na execução.\\
Distância selecionada & \num{6112} km / \num{3300.216} nm & Distância usada para calcular a navegação.\\
\bottomrule
\end{tabularx}
\end{table}

\section{Comparação com uma referência externa}
\label{sec:benchmark-externo}

O sistema também compara seus resultados com uma planilha produzida fora do projeto. Essa comparação externa recebe o nome de \emph{benchmark}. O workbook associado aos estudos do Prof. Dr. Gustavo Costa permite verificar se as duas análises apontam o mesmo modo como menos emissor e se os valores estão na mesma ordem de grandeza \cite{competitiveness2024,decarb2024}. Ele não fornece uma verdade de referência para calibrar o sistema.

A matriz possui \num{36} células. Seis comparam um porto com ele mesmo. Outras nove têm valor rodoviário nulo ou não positivo. Esses 15 casos não permitem comparação numérica. Restam \num{21} pares OD positivos e suportados.

Nos 21 casos, as duas fontes apontam a mesma direção para emissões: a alternativa com cabotagem apresenta menor valor. Entretanto, existe diferença relevante de magnitude. A classificação \codecell{same\_direction\_large\_gap} significa exatamente isso: o sinal é coerente, mas os números não são equivalentes.

Para comparar magnitudes seria necessário alinhar carga, distância, portos, veículo rodoviário, intensidade marítima, alocação e operações portuárias. Também seria necessário alinhar o que cada cálculo inclui e exclui. Esse limite recebe o nome de \emph{fronteira}. Sem esse alinhamento, o benchmark informa apenas a direção do resultado. Seus parâmetros não substituem os parâmetros do \CabotageLens{}.

\section{Categorias de uso da evidência}
\label{sec:categorias-uso}

Cada resultado recebe uma categoria que limita o que pode ser afirmado. A categoria funciona como uma instrução de leitura. Os termos em inglês da Tabela~\ref{tab:categorias-uso} são códigos internos gravados junto ao resultado. Eles não representam novas fórmulas. Essa classificação evita que um valor disponível seja apresentado como validação universal.

\begingroup
\small
\setlength{\LTleft}{0pt}
\setlength{\LTright}{0pt}
\setlength{\tabcolsep}{4pt}
\renewcommand{\arraystretch}{1.15}
\begin{longtable}{L{0.29\textwidth}L{0.32\textwidth}L{0.31\textwidth}}
\caption{Categorias de uso e limites de interpretação.}
\label{tab:categorias-uso}\\
\toprule
Categoria & Uso permitido & Uso proibido \\
\midrule
\endfirsthead
\toprule
Categoria & Uso permitido & Uso proibido \\
\midrule
\endhead
\codecell{headline\_candidate} & Resultado principal somente quando houver evidência de rota, fonte e fronteira suficiente. & Promover automaticamente todo cenário executado.\\
\codecell{sensitivity\_discussion} & Discutir comportamento sob uma premissa identificada. & Tratar a premissa como linha de base ou calibração.\\
\codecell{limitation\_example} & Explicar limites, como caso de origem e destino no mesmo porto. & Usar como desempenho normal da cabotagem.\\
\codecell{excluded} & Registrar que o cenário está fora da fronteira definida. & Produzir conclusão numérica para o cenário excluído.\\
\codecell{reference\_needed} & Identificar a evidência necessária para uso quantitativo. & Substituir a referência ausente por porto ou fonte diferente sem rotulagem.\\
\codecell{methodology\_blocked} & Registrar uma decisão metodológica ainda não resolvida. & Converter o caso em resultado sem regra rastreada.\\
\codecell{benchmark\_supports\_direction} & Sustentar consistência direcional cautelosa. & Alegar reprodução exata ou validade operacional.\\
\codecell{benchmark\_boundary\_mismatch} & Preservar diferenças entre \ttw{}, \wtw{}, LCA, CO$_2$ e \cooe{}. & Tratar fronteiras ou gases distintos como equivalentes.\\
\codecell{not\_comparable} & Justificar a ausência de comparação quantitativa. & Inferir desempenho a partir de valores incompatíveis.\\
\bottomrule
\end{longtable}
\endgroup

Os testes mostram se o código segue as regras. A auditoria mostra quais dados sustentam o número. O benchmark mostra apenas a direção em relação a uma fonte externa. Nenhuma dessas evidências, isoladamente, comprova superioridade universal da cabotagem, disponibilidade de serviço ou equivalência com tarifas comerciais.
